% Generated draft: requires mathematical and source review.
\documentclass[11pt]{article}
\usepackage[margin=25mm]{geometry}
\usepackage{fontspec}
\setmainfont{Latin Modern Roman}
\usepackage{amsmath,amssymb,mathtools,unicode-math}
\setmathfont{Latin Modern Math}
\usepackage{enumitem,needspace,xurl,hyperref}
\setlength{\parindent}{0pt}
\setlength{\parskip}{5pt}
\setlength{\emergencystretch}{3em}
\newcommand{\N}{\mathbb{N}}
\newcommand{\Z}{\mathbb{Z}}
\newcommand{\Q}{\mathbb{Q}}
\newcommand{\R}{\mathbb{R}}
\newcommand{\C}{\mathbb{C}}
\newcommand{\F}{\mathbb{F}}
\newcommand{\A}{\mathbb{A}}
\newcommand{\PP}{\mathbb{P}}
\newcommand{\E}{\mathbb{E}}
\newcommand{\eps}{\varepsilon}
\newcommand{\dd}{\,\mathrm{d}}
\begin{document}
\section*{UnsolvedMath ID: 701757}
\textbf{Source problem code:} catalogue entry
\subsection{Definitions and mathematical statement}
Let \(G\) be a group and let \(k\in\Z\) with \(k>1\). For each \(i\in\{1,\ldots,k\}\), let \(G_i\leq G\) be a subgroup of finite index \(n_i=[G:G_i]\), where the index is the number of left cosets of \(G_i\) in \(G\), and let \(a_i\in G\). Define \(a_iG_i=\{a_ig:g\in G_i\}\). Suppose these left cosets are pairwise disjoint:
\[
(a_iG_i)\cap(a_jG_j)=\varnothing\qquad(1\leq i<j\leq k).
\]
The conjecture of Z.-W. Sun asserts that, for every such choice, there exist indices \(1\leq i<j\leq k\) such that
\[
\gcd(n_i,n_j)\geq k,
\]
where \(\gcd\) denotes the greatest common divisor. No finiteness assumption on \(G\), normality assumption on the subgroups, or requirement that the cosets cover \(G\) is imposed. The supplied source states that the conjecture holds for \(k=2,3,4\); this reported qualification is not independently verified here.
\subsection{Short English statement}
Z.-W. Sun conjectures that among any \(k>1\) pairwise disjoint left cosets of finite-index subgroups of a group, two subgroup indices have greatest common divisor at least \(k\). The supplied record labels the general conjecture open and reports that it holds for \(k=2,3,4\).
\subsection{Sources}
\begin{itemize}
\item Local UnsolvedMath registry, supplied record with permanent ID \(701757\), exact source problem code \texttt{catalogue entry}, titled \emph{Kourovka Notebook Problem 20.99}. The record supplies the status \texttt{open} and attributes the problem to Z.-W. Sun, proposed in 2022.
\item \emph{The Kourovka Notebook (Unsolved Problems in Group Theory)}, Problem 20.99, attributed to Z.-W. Sun. The supplied background identifies the 20th issue (2022) and the 21st issue edition (2026), with identifier \texttt{arXiv:1401.0300v48}. The supplied collection description names E. I. Khukhro and V. D. Mazurov as editors. These bibliographic details and the reported cases \(k=2,3,4\) are supplied source information, not independently verified here.
\end{itemize}
\end{document}
